\section{Flagship Receipts}
\sectionkicker{FLAGSHIP BOUNDS}
\begin{wideflow}
{\Large\bfseries 看板の主張と受け止め方}\par
\smallskip
{\color{SurveyMuted}旗艦の看板と道具の中の配備を並べ、測定値の出どころと留保を切り離さない。}\par
\end{wideflow}

8月24日に報じられたClaude Fable 5は、評価機関がMythos級の初号として位置づけ、GDPval-AAの道具的知識労働で首位としたことが伝えられる\cite{fable5}。Mythos級は評価機関の枠組みであり業界標準ではない。同じ評価機関のAnalystAgentではGemini 3.7 Flashが高く、Opus 5、Fable 5と続くという内訳が示され、Fable 5の強みは開かれた道具的推論にあり、表計算の操作ではないという留保が付く。いずれも一次資料の確認を待つ報道経由の主張であり、同週に並んだ他モデルの一覧は背景の密度として触れるだけで個別の扱いにはしない。

\begin{claimboundary}[未解決の境界]
性能や等級の主張は評価機関と報道を経由したものであり、一次資料の確認を待つ。Mythos級を業界標準として扱わず、Terminal-Benchの手法を作り込まない。
\end{claimboundary}

同じ24日に一次発表で確かめられるのは、GPT-5.6の一族がAWSのKiroで使えるようになったことである\cite{gpt56kiro}。要件から設計、仕事、性質検査、区切りの審査へ進む仕様駆動の道具という位置づけも一次資料で確認できる。Kiro上でのTerminal-Bench 2.1の仕事をGPT-5.6 Terraが約82\%の費用減で終えたという数値は、提供側の共同測定としての主張であり、手法の公開や独立の再現を欠く。価格や地域、模型票の事実はこの発表では定まらない。
